\subsection{de Sitter Vacua in the String Landscape --- arXiv:1908.05288}

\textbf{Authors:} Keshav Dasgupta, Maxim Emelin, Mir Mehedi Faruk, Radu Tatar

\paragraph{主な主張}
内部空間とフラックスの時間依存性および量子効果を含めたtype IIB/M理論の解析から、Newton定数を一定に保ちつつ四次元de Sitter等長性を実現する可能性を論じる\cite{Dasgupta:2019gcd}。

\paragraph{新規性}
静的な内部空間で生じる量子補正の階層性の問題を拡張し、摂動・非摂動および局所・非局所補正を含む時間依存背景を検討する。

\paragraph{位置付け}
de Sitter構成の可否を、量子補正の制御と有効場の理論の成立条件から検討するtop-down研究。

\paragraph{コメント（読書メモ）}
当時の読書メモでは、de Sitter真空の明示的なtop-down構成の難しさと、弦ループ補正が有効な宇宙項へ時間依存の寄与を与える点に注目した。単純な設定に基づくde Sitter予想が、KKLTのanti-D3-braneなどの非自明な効果を含む場合にも適用できるかを検討し、複雑な量子補正を調べる必要があると考えた。先行解析\cite{Dasgupta:2018rtp}では、考察した時間非依存の設定に無限個の補正が同程度に寄与し、有限次数で制御した有効理論の記述が難しくなる。本論文ではdipole型・Kasner型変形だけではこの問題が改善しない一方、内部空間やフラックスの適切な時間依存性により、各次数に寄与する補正の階層性を回復できる可能性を検討している、と理解した。
